\documentclass[12pt, a4paper, titlepage]{article}

\usepackage{hyperref}

\title{
  \Huge{INFO 250 - 002 Web Development I Final Project Proposal (Spring 2026)}\\
\huge{\textbf{Quizzing Website}}}
\author{\textbf{Group 3:}\\
 Singyang HOEURNG\\
 Kimchour LUY\\ 
 Sopheakchesda LY\\
 Penhsang TAO\\
}

\date{\today}

\begin{document}

\maketitle

\section*{About}

We'll make a simple quizzing website with support for \texttt{html} elements such as audio, images, and embedded elements. We will only support multiple choice questions and perhaps text-based questions for now.

\section*{Abstract}

This site will be built with pure \texttt{html}, pure \texttt{css} and very minimal \texttt{javascript}. There will be a minimum of 4 pages containing the following content:

\begin{enumerate}
  \item Home: there will be a hero section, a short description, and a list of quizzes available
  \item Documentation: contains the correct \texttt{html} format for quiz creation, such that it's easy and repeatable to create new quizzes (We'll simply inject that \texttt{html} into the body of a simple template page we define for now).
\item Submit a quiz: we'll have a form to allow visitors to submit their quizzes to be hosted on our website
  \item \texttt{mp3} vs \texttt{wav}: a quiz to test whether the test taker can hear the difference between the lossy \texttt{mp3} audio vs the lossless \texttt{wav} audio format. We may also include different lossy audio compression formats such as \texttt{m4a}, \texttt{opus}, \texttt{ogg}, with different bitrates. This quiz demonstrates the usage of audio elements.
  \item Guess the song quiz
  \item \texttt{png} vs \texttt{jpg}: a quiz to test whether we can see the difference between the lossless \texttt{png} image compression format vs the \texttt{jpg} image compression format. We may also include the \texttt{heic} and use different compression ratios
  \item Guess the flag quiz
  \end{enumerate}

 We'll make sure to: 

  \begin{itemize}
    \item format the \texttt{html} page with proper semantic elements
    \item use external \texttt{css} for styling 
    \item make the site mobile friendly through responsive layouts like \texttt{flexbox} and \texttt{css}
    \item use fonts using Google Fonts
    \item include tables in the quizzes
    \item include images and multimedia elements in the quizzes
    \item include forms to report issues, and to submit quizzes will be used
    \item include proper \texttt{css} styling will be used to make the site easy to use
  \end{itemize}


  To handle the logic, we'll use some \texttt{javascript}, but only to do some basic checking of whether the user gave the right answer, and to display the final score. 


\section*{References}

No references other than the \href{https://aupp.instructure.com/courses/4360/files/621169?wrap=1}{sample project proposal} and the \href{https://aupp.instructure.com/courses/4360/files/610452?wrap=1}{final project guidelines} were used.

\section*{AI usage}

No AI was used to make this proposal.

\end{document}
